\documentclass{article}
\usepackage{thumbpdf,lmodern}
\usepackage{amsmath}
\usepackage{graphicx}
\usepackage[margin=1in]{geometry}

\newcommand{\pkg}[1]{\texttt{#1}}
\newcommand{\proglang}[1]{\textsf{#1}}
\newcommand{\code}[1]{\texttt{#1}}
\newcommand{\Plainauthor}[1]{}
\newcommand{\Plaintitle}[1]{}
\newcommand{\Shorttitle}[1]{}
\newcommand{\Plainkeywords}[1]{}
\newcommand{\Address}[1]{}
\newcommand{\Email}[1]{\texttt{#1}}
\newcommand{\myabstract}{}
\newcommand{\Abstract}[1]{\renewcommand{\myabstract}{#1}}
\newcommand{\mykeywords}{}
\newcommand{\Keywords}[1]{\renewcommand{\mykeywords}{#1}}
\newcommand{\Title}[1]{\title{#1}}

\author{Dylan A. Mordaunt\\Victoria University of Wellington}
\Title{Supplement A: Mathematical Proofs and Diagnostics}
\Abstract{
  This supplement provides the detailed mathematical derivation of the Value of Perspective (VoP) metric and presents convergence diagnostics for the probabilistic sensitivity analysis (PSA) simulations.
}
\Keywords{VoP, convergence, diagnostics}

\begin{document}
\maketitle
\begin{abstract}
\myabstract
\end{abstract}

\section{Mathematical Derivation of Value of Perspective}

The Value of Perspective (VoP) is derived from the principles of Value of Information (VOI) analysis.
Let $\theta$ be the vector of uncertain parameters.
Let $d \in D$ be the set of possible decisions (e.g., fund, do not fund).
Let $NB(d, \theta, P)$ be the net benefit of decision $d$ given parameters $\theta$ under perspective $P$.

We assume the societal perspective ($P_{soc}$) represents the true social welfare.
The optimal decision under the societal perspective, $d^*_{soc}(\theta)$, maximizes the expected net benefit:
\begin{equation}
d^*_{soc} = \arg\max_{d} E_{\theta}[NB(d, \theta, P_{soc})]
\end{equation}

However, decision-makers often operate under a restricted health system perspective ($P_{hs}$).
The decision chosen under this perspective, $d^*_{hs}$, is:
\begin{equation}
d^*_{hs} = \arg\max_{d} E_{\theta}[NB(d, \theta, P_{hs})]
\end{equation}

The \emph{Value of Perspective} is the expected opportunity loss incurred by making the decision $d^*_{hs}$ instead of $d^*_{soc}$, evaluated from the societal perspective:
\begin{equation}
VoP = E_{\theta}[NB(d^*_{soc}, \theta, P_{soc})] - E_{\theta}[NB(d^*_{hs}, \theta, P_{soc})]
\end{equation}

If $d^*_{soc} = d^*_{hs}$, the VoP is zero. If the perspectives lead to different decisions, the VoP quantifies the societal benefit foregone.

\section{Convergence Diagnostics}

To ensure the stability of the VoP and other probabilistic metrics, we monitored the convergence of the Incremental Net Monetary Benefit (INMB) over the 1,000 PSA iterations.
Figure~\ref{fig:convergence} displays the running mean and 95\% confidence intervals for the INMB of the \emph{HPV Vaccination} case study (Societal Perspective).

\begin{figure}[ht]
\centering
\includegraphics[width=0.8\textwidth]{figures/convergence_hpv_vaccination.png}
\caption{Convergence of Incremental Net Monetary Benefit (Societal Perspective) for HPV Vaccination. The running mean stabilizes after approximately 500 iterations.}
\label{fig:convergence}
\end{figure}

Similar convergence was observed for all other interventions (plots available upon request).

\end{document}
